\newpage
\section*{РЕФЕРАТ}

ВКР \pageref{LastPage} с., \total{figure} рис., \total{equation} ф-л, \total{table} табл., \total{bibitems}
источн., 1 прил.

КЛЮЧЕВОЕ СЛОВО 1, КЛЮЧЕВОЕ СЛОВО 2, КЛЮЧЕВОЕ СЛОВО 3, КЛЮЧЕВОЕ СЛОВО 4, КЛЮЧЕВОЕ СЛОВО 5, КЛЮЧЕВОЕ
СЛОВО 6, КЛЮЧЕВОЕ СЛОВО 7, КЛЮЧЕВОЕ СЛОВО 8

Объектом исследования является процесс или явление, рассматриваемое в рамках темы выпускной квалификационной работы.

Цель работы – разработка и обоснование предлагаемого решения поставленной научно-технической задачи.

В процессе работы проводились: анализ предметной области, обзор существующих подходов, разработка метода или алгоритма,
программная реализация и оценка полученных результатов.

Научная ценность работы заключается в предложении нового или усовершенствованного подхода к решению рассматриваемой
задачи.

Основные конструктивные и технико-эксплуатационные показатели: перечислите измеримые характеристики разработанного
решения с указанием единиц измерения.

Степень внедрения – опишите результаты апробации, пилотного внедрения или экспериментальной проверки.

Эффективность определяется количественными показателями качества, например точностью, временем обработки или
экономическим эффектом.
